\SourceLeaf{213}%
\noindent setzen. Wenn man hier die allgemeine Auflösung des 5. Artikels zur Anwendung bringt, so findet man, dass, indem \(f\) eine willkürliche Function bedeutet, \(T\) dem reellen und \(i\log\mathrm{cotang}\tfrac{1}{2}U\) dem imaginären Theile von
\[
f(t+i\log\{\mathrm{cotang}\tfrac{1}{2}w\mathord{.}(\frac{1-\varepsilon\cos w}{1+\varepsilon\cos w})^{\frac{1}{2}\varepsilon}\})
\]
\noindent gleich gesetzt werden muss\footnote{Wir übergehen hier theils die zweite Auflösung des 5. Artikels, die sich von der obigen nur durch Vertauschung von \(-T\) gegen \(+T\) unterscheiden und einer verkehrten Darstellung entsprechen würde, theils den Fall eines länglichen Ellipsoids, dessen Behandlung nach dem, was im vorigen Art. vorgekommen, sich aus der des abgeplatteten von selbst ergibt.}.

Die einfachste Auflösung wird sein, \(f\upsilon=\upsilon\) zu setzen, wodurch
\[
T=t,\qquad \mathrm{tang}\tfrac{1}{2}U=\mathrm{tang}\tfrac{1}{2}w\mathord{.}(\frac{1+\varepsilon\cos w}{1-\varepsilon\cos w})^{\frac{1}{2}\varepsilon}
\]
\noindent wird. Dies bietet eine für die höhere Geodäsie überaus brauchbare Transformation dar, von welcher Benutzung wir jedoch hier nur einiges und nur kurz andeuten können. Wenn nemlich auf der Oberfläche des Ellipsoids und der Kugel diejenigen Punkte als einander correspondirend angesehen werden, die einerlei Länge haben, und deren Breiten resp. \(90^0-w\), \(90^0-U\), vermöge der angeführten Gleichung zusammenhangen, so entspricht einem System von, verhältnissmässig, kleinen Dreiecken (und das werden diejenigen immer sein, die zur wirklichen Messung dienen können), die auf der Oberfläche des Sphäroids durch kürzeste Linien gebildet werden, auf der Kugelfläche ein System von Dreiecken, deren Winkel den correspondirenden auf dem Sphäroid \textit{genau} gleich sind, und deren Seiten von grössten Kreisbogen so wenig abweichen, dass sie in den meisten Fällen, wo nicht die alleräusserste Schärfe verlangt wird, als damit zusammenfallend betrachtet werden können, so wie auch da, wo die grösste Genauigkeit gefordert wird, die Abweichung vom grössten Kreise leicht mit aller nöthigen Schärfe durch einfache Formeln sich berechnen lässt. Man kann daher das ganze System, nachdem man zuerst eine Dreiecksseite auf die Kugelfläche gehörig übertragen hat, ganz so, als wenn es auf dieser selbst läge, vermittelst der Winkel berechnen, nöthigenfalls mit der eben angedeuteten Modification, für alle Punkte des Systems die Werthe von \(T\) und \(U\) bestimmen, und von letztern auf die correspondirenden Werthe von \(w\) (am einfachsten vermittelst einer äusserst leicht zu construirenden Hülfstafel) zurückgehen.
